\subsection{A specified localization at rank four}\label{sec:rank4-word-minor}

We now give the finite data for one localization. The generic dimension is a published input: \citet[\S11.8, pp.~67--68]{BSM2013} gives $\dim_{K_v}A_{K_v}=114$, with simple-module dimensions
\[
 1,2,3,6,5,5,2,3,1.
\]
Their squared dimensions sum to $114$. The calculation below independently proves that the displayed words are independent; the published result supplies the upper bound and hence their spanning property.

List the permutations of $(1,2,3,4)$ lexicographically as $\pi_0,\ldots,\pi_{23}$. Number the standard tensor coordinates by
\[
 24i+j\quad\longleftrightarrow\quad T_{\pi_i}\otimes T_{\pi_j},
 \qquad 0\le i,j<24.
\]
For a string $a_1\cdots a_k$ in $\{1,2,3\}$, the corresponding word is $P_{s_{a_1}}\cdots P_{s_{a_k}}$; the empty string denotes $1$. Table~\ref{tab:rank4-word-minor} specifies $114$ ordered pairs $(w_i,r_i)$. Let
\[
 C(v)=\bigl([w_j]_{r_i}\bigr)_{1\le i,j\le114},\qquad
 g(v)=\det C(v),
\]
where $[w]_{r}$ is the indicated coordinate of the word. Define the Laurent expression
\begin{equation}\label{eq:rank4-norm-minor}
 f(D)=g(v)g(-v)g(v^{-1})g(-v^{-1}),\qquad
 D=v^2+2+v^{-2}.
\end{equation}
Thus the complete finite definition of $f$ is the Hecke multiplication rule together with the table; no expanded polynomial is required.

\begin{proposition}\label{prop:rank4-word-minor}
The expression in \eqref{eq:rank4-norm-minor} is a nonzero polynomial in $\mathbb Z[D]$. The words of Table~\ref{tab:rank4-word-minor} form an $F_D$-basis of $A_{F_D}$, and
\[
 A_{S_D}[1/f]=\bigoplus_{i=1}^{114}S_D[1/f]w_i.
\]
In particular $S_D[1/f]$ is an explicitly specified coefficient ring for this word basis at rank four. This assertion does not require that $f$ be a least denominator or that the word basis be positive.
\end{proposition}

\begin{proof}
All entries of $C(v)$ lie in $R_v$. Exact arithmetic modulo the prime $p=1000003$ gives
\[
\begin{array}{c|rrrr}
v&2&-2&2^{-1}&-2^{-1}\\\hline
g(v)\bmod p&274637&274637&274637&274637
\end{array}
\]
and consequently
\[
 f(25/4)\equiv988208\pmod p.
\]
The rational number $25/4$ is evaluated using the inverse of $4$ modulo $p$. These nonzero residues prove $g\ne0$ and $f\ne0$ as Laurent polynomials. The $114$ word columns are therefore independent over $K_v$. The published upper bound $114$ proves that they are a $K_v$-basis; Theorem~\ref{thm:hecke-word-descent} then makes them an $F_D$-basis.

For completeness, the modular calculation is a finite certificate with no floating-point comparison. On the standard Hecke basis, right multiplication by $T_s$ sends $T_w$ to $T_{ws}$ when length increases, and to $T_{ws}+(v-v^{-1})T_w$ when length decreases. Expanding the two tensor summands defining $P_s$ computes each word column. Restrict to the listed rows and perform elimination in $\mathbb F_p$; the four displayed values are the resulting determinants. The program \emph{Rank-four word coordinates} records the words, rows and all four values. A separate program, \emph{Rank-four minor replay}, reconstructs the same columns by left multiplication in reversed word order, expands the two tensor summands separately, and uses modular Bareiss elimination. It gives the same four determinants and norm. Both programs and their exact output certificates are included in the accompanying workbench. Trial division up to $\lfloor\sqrt p\rfloor$ also checks primality.

The four factors of $f$ are permuted by $v\mapsto-v$ and $v\mapsto v^{-1}$. Hence $f$ lies in the invariant subring $R_v^{\operatorname{Gal}(K_v/F_D)}$. This subring is exactly $S_D$: sign invariance removes all odd powers of $v$, and inversion invariance writes the result as an integer linear combination of $1$ and $v^{2n}+v^{-2n}$. If $z=v^2+v^{-2}=D-2$, these latter expressions are polynomials in $z$, by the recurrence
\[
 a_0=2,\qquad a_1=z,\qquad a_{n+1}=za_n-a_{n-1}.
\]
Thus $f\in\mathbb Z[D]$, and also $R_v\cap F_D=S_D$.

Let $w$ be any word and write $w=\sum_jx_jw_j$ with $x_j\in F_D$. Apply Cramer's rule to its selected coordinate equations over $R_v$. It gives $g(v)x_j\in R_v$. The quotient $f/g$ is the product of the other three conjugate minors and lies in $R_v$. Therefore $fx_j\in R_v\cap F_D=S_D$. Every word consequently belongs to $f^{-1}\sum_jS_Dw_j$. Localizing gives one inclusion in the asserted equality; the reverse inclusion follows because every $w_j$ already belongs to $A_{S_D}$. Independence was proved above.
\end{proof}

\begin{table}[!ht]
\centering\footnotesize
\setlength{\tabcolsep}{3pt}
\begin{tabular}{lr lr lr lr lr lr }
\toprule
word & row & word & row & word & row & word & row & word & row & word & row\\
\midrule
$\varnothing$ & 0 & \texttt{1} & 6 & \texttt{2} & 2 & \texttt{3} & 1 & \texttt{12} & 8 & \texttt{13} & 7\\
\texttt{21} & 12 & \texttt{23} & 3 & \texttt{32} & 4 & \texttt{121} & 14 & \texttt{123} & 9 & \texttt{132} & 10\\
\texttt{212} & 50 & \texttt{213} & 13 & \texttt{232} & 5 & \texttt{321} & 18 & \texttt{323} & 25 & \texttt{1212} & 54\\
\texttt{1213} & 15 & \texttt{1232} & 11 & \texttt{1321} & 20 & \texttt{1323} & 31 & \texttt{2121} & 56 & \texttt{2123} & 51\\
\texttt{2132} & 16 & \texttt{2321} & 19 & \texttt{2323} & 26 & \texttt{3212} & 52 & \texttt{3213} & 33 & \texttt{3232} & 27\\
\texttt{12121} & 60 & \texttt{12123} & 55 & \texttt{12132} & 17 & \texttt{12321} & 21 & \texttt{12323} & 32 & \texttt{13212} & 61\\
\texttt{13213} & 28 & \texttt{13232} & 34 & \texttt{21213} & 57 & \texttt{21232} & 30 & \texttt{21321} & 22 & \texttt{21323} & 35\\
\texttt{23212} & 36 & \texttt{23213} & 37 & \texttt{23232} & 29 & \texttt{32121} & 58 & \texttt{32123} & 53 & \texttt{32132} & 38\\
\texttt{32321} & 59 & \texttt{121213} & 63 & \texttt{121232} & 39 & \texttt{121321} & 23 & \texttt{121323} & 40 & \texttt{123212} & 41\\
\texttt{123213} & 42 & \texttt{123232} & 43 & \texttt{132121} & 66 & \texttt{132123} & 62 & \texttt{132132} & 44 & \texttt{132321} & 64\\
\texttt{212132} & 45 & \texttt{212321} & 65 & \texttt{212323} & 67 & \texttt{213212} & 75 & \texttt{213213} & 76 & \texttt{213232} & 77\\
\texttt{232121} & 78 & \texttt{232123} & 79 & \texttt{232132} & 80 & \texttt{232321} & 100 & \texttt{321213} & 81 & \texttt{321232} & 82\\
\texttt{321321} & 84 & \texttt{321323} & 101 & \texttt{323212} & 152 & \texttt{1212132} & 46 & \texttt{1212321} & 69 & \texttt{1212323} & 68\\
\texttt{1213212} & 83 & \texttt{1213213} & 85 & \texttt{1213232} & 86 & \texttt{1232121} & 87 & \texttt{1232123} & 88 & \texttt{1232132} & 89\\
\texttt{1232321} & 90 & \texttt{1321213} & 91 & \texttt{1321232} & 92 & \texttt{1321321} & 102 & \texttt{1321323} & 103 & \texttt{1323212} & 104\\
\texttt{2121321} & 105 & \texttt{2121323} & 106 & \texttt{2123212} & 107 & \texttt{2123213} & 150 & \texttt{2123232} & 151 & \texttt{2132121} & 108\\
\texttt{2132123} & 153 & \texttt{2132132} & 109 & \texttt{2132321} & 126 & \texttt{2321232} & 154 & \texttt{2321321} & 155 & \texttt{12121321} & 110\\
\texttt{12121323} & 111 & \texttt{12123212} & 112 & \texttt{12123213} & 113 & \texttt{12123232} & 156 & \texttt{12132121} & 157 & \texttt{12132123} & 114\\
\texttt{12132132} & 127 & \texttt{12132321} & 128 & \texttt{12321232} & 158 & \texttt{12321321} & 159 & \texttt{21213212} & 160 & \texttt{121213212} & 161\\
\bottomrule
\end{tabular}
\caption{Ordered word--coordinate pairs defining the rank-four minor. Read the pairs from left to right across each row, then from top to bottom. The empty word is the unit; coordinate indices begin at zero.}
\label{tab:rank4-word-minor}
\end{table}


This certificate supplies one finite localization for the specified integral word order. It leaves open whether a smaller localization suffices, whether the order is free before localization, and whether a distinguished positive basis exists over an appropriate coefficient ring. None of those questions follows from modular nonvanishing or from the published generic dimension.
\clearpage
